\section{引言}
\label{sec:introduction}
单风场功率预测的输出是某个目标场站未来的发电轨迹，但这一输出定义并不限定模型必须只用该站的数据训练。跨风场共享参数可以扩大训练样本，而部署时仍可能只需目标风场的历史功率、本站天气预报与静态信息。与之不同，若模型在线读取区域聚合功率，则会增加对其他风场数据可用性的依赖。明确这两种信息来源，有助于判断预测收益能否转化为实际的单站部署能力。

历史功率与数值天气预报（NWP）提供互补的信息。前者反映已经实现的运行水平及短时变化，后者描述未来气象条件。在逐小时天气预测未来六小时、每十五分钟功率的任务中，插值可以对齐时间，却不能说明天气应如何改变已有的功率判断。本文关注的核心问题是：当历史模型已给出一条未来参考轨迹时，将这条轨迹显式传递给天气分支，是否有助于形成更准确的多步修正？

为此，本文先训练历史预测器，再冻结其参数。修正网络在每个提前量读取参考功率、相邻参考差分、提前量及历史记忆均值，联合对齐后的天气特征形成条件调整；解码器同时读取完整历史记忆。参考轨迹既是修正条件，又是最终输出的加性基准。这一组织方式将历史支持的未来判断显式暴露给天气分支。它没有增加独立观测，也不将历史模型解释为无天气影响的物理反事实；修正项允许同时承担天气相关预测和基准校准。

本文以共享训练、本站输入的 CBH-SL 为单站部署的主要实证对象，并用独立训练及区域历史输入构成四格比较。进一步在保留区域历史的版本上检验条件入口与非线性天气编码。这些实验回答不同问题：本地部署表现由 CBH-SL 支撑，模块归因由各自匹配的对照支撑，区域增强版本的优势不能直接转移为尚未评估的本站输入增强模型。类似地，源站误差更低也不意味着在新站上更容易适配，因此未见风场实验同时保留强预训练参照与负迁移实例。

本文的贡献概括如下：
\begin{itemize}
 \item \textbf{轨迹条件化的天气信息组织。}在保留完整历史解码记忆的同时，将历史模型的未来轨迹及其邻步差分显式用于天气修正。匹配条件对照识别这一信息包的收益，并将其与较小的摘要增量及单纯输出相加增量区分开。
 \item \textbf{训练共享与在线输入的分离验证。}以单风场预测为目标，通过共享/独立训练和本站/区域历史的四格设计，量化参数共享与在线区域输入各自对应的效果，为本站部署给出可复核的证据。普通参数共享本身不作为独立算法创新。
 \item \textbf{跨风场适配的收益及边界。}在不同适配训练预算下评价 20 个未见风场，并结合不依赖模型误差的代表站选择、完整场站统计和失败案例，揭示相对从头训练的收益与相对强预训练模型的不足。
\end{itemize}
主结论使用全部源站或最终目标站，不由少数典型站重新定义评价集合。天气编码的非线性、分时段表现与适配结果用于检验上述主张，不额外包装为未经独立验证的物理机制。
